\ExamPart{Trắc nghiệm trả lời ngắn}{2}{Thí sinh trả lời từ câu 1 đến câu 4.}

\NQuestion Trong mặt phẳng tọa độ $Oxy$ cho hai đường thẳng
\[
\Delta_1:2x-y+1=0
\quad \text{và} \quad
\Delta_2:x+2y-4=0.
\]
Gọi $\alpha$ là góc giữa hai đường thẳng $\Delta_1$ và $\Delta_2$. Tìm $\tan\alpha$.

\NQuestion Từ các chữ số $0,1,2,3,4,5,6$. Hỏi có thể lập được bao nhiêu số tự nhiên lẻ có hai chữ số khác nhau?

\NQuestion Cho phương trình
\[
\sqrt{x^2-4x+1}=\sqrt{x+2}
\]
có hai nghiệm $x_1$ và $x_2$, với $x_1<x_2$. Tính
\[
2x_1+3x_2.
\]

\NQuestion Cho bất phương trình
\[
\left(x^2+2\right)\left(-2x^2+3x+2\right)\ge 0
\]
có tập nghiệm
\[
S=\left[a;b\right].
\]
Tính
\[
2a+5b.
\]

\ExamPart{Tự luận}{3}{Thí sinh trả lời từ câu 1 đến câu 3.}

\NQuestion Giải bất phương trình sau:
\[
\dfrac{x^2-4x+3}{x-4}>0.
\]

\NQuestion Trong mặt phẳng tọa độ $Oxy$, cho hình chữ nhật $ABCD$ có
\[
A\left(-2;1\right),\quad B\left(2;-1\right),\quad C\left(0;-5\right).
\]
\begin{SubQuestions}
\item Tìm tọa độ giao điểm $I$ của hai đường chéo $AC$ và $BD$.
\item Tính chu vi hình chữ nhật $ABCD$.
\end{SubQuestions}

\NQuestion Trong mặt phẳng tọa độ $Oxy$, cho
\[
A\left(2;1\right),
\]
và đường thẳng có phương trình
\[
d:2x+5y-9=0.
\]
\begin{SubQuestions}
\item Viết phương trình đường thẳng $\Delta$ đi qua điểm $A$ và vuông góc với đường thẳng $d$.
\item Trên đường thẳng $d$ lấy $M(2;1)$. Gọi $H$ là giao điểm của đường thẳng $\Delta$ và đường thẳng $d$. Tính diện tích tam giác $AHM$.
\end{SubQuestions}
